% Literature-only draft: no round-2 outcomes or empirical literature numbers.
% Primary-source verification is recorded beside the relevant bibliography entries.
\section{Related Work}
\label{sec:related}

\subsection{Threshold-aware evaluation in speech and biometrics}
% src: bib:bengio04_hter
The distinction between selecting a decision rule and evaluating it is established in biometric authentication.
\citet{bengio04_hter} evaluate HTER on a test set after selecting an EER threshold on separate validation data, and
show why uncertainty in the class-specific errors matters for statistical comparison. Face presentation-attack
detection also uses development-set thresholds followed by test-set HTER \citep{chingovska12_lbp}.
For speech, \citet{korshunov16_crossdb} use the same separation when transferring countermeasures between databases
and between logical-access and replay attacks. These precedents motivate keeping the threshold-selection data
explicit when assessing a speech deepfake detector.

% src: bib:brummer06_calibration
Speaker-recognition evaluation further distinguishes discrimination, described by the attainable error trade-offs,
from calibration, which determines how well a score supports a particular decision rule
\citep{brummer06_calibration,vanleeuwen07_intro}. The difference between actual cost and minimum cost expresses this
distinction for a specified application. A minimum found using evaluation labels describes what an oracle threshold
could achieve; actual cost evaluates the supplied scores at the application's threshold. Our TTG follows this
lineage at the equal-cost operating point: it measures the excess HTER of a development threshold above the target's
minimum HTER. It is a threshold-specific loss, rather than a claim that the CM scores are calibrated likelihood
ratios across applications.

% src: bib:wang20_asvspoof2019
The ASVspoof challenges provide a complementary evaluation lineage. ASVspoof~2019 and ASVspoof~2021 report EER and
minimum tandem detection cost, with t-DCF assessing a countermeasure together with an ASV system; the 2021 deepfake
(DF) track, which involves no ASV system, reports EER only \citep{kinnunen18_tdcf,wang20_asvspoof2019,liu23_asvspoof2021}. The stand-alone detection track of ASVspoof~5 reports
minimum DCF, actual DCF, and log-likelihood-ratio cost \citep{wang24_asvspoof5}. Its actual DCF uses a fixed Bayes
threshold under an LLR interpretation of the submitted scores. This differs from transferring an empirical
development threshold, but both protocols expose behavior that a threshold optimized on the target cannot reveal.

% src: bib:zhou26_eer_hides
A direct speech-deepfake precedent is \citet{zhou26_eer_hides}, who freeze an SSL-AASIST detector, select a source
EER threshold on labelled ASVspoof~2019 LA evaluation data, and transfer it to In-the-Wild and a subset of
ASVspoof~2021 DF. They report HTER and its false-rejection/false-acceptance decomposition, and audit unlabelled
score corrections. Their monotone-invariance argument clarifies that a strictly increasing score transform can
alter decisions at a fixed numerical threshold while leaving the ROC and EER unchanged. Our audit uses development data for
threshold selection and extends the assessment across checkpoints, policies, training runs, and controlled
perturbations. The contribution is the scope and registration of this assessment, rather than the invention of
transferred-threshold evaluation.

\subsection{Generalization, dataset cues, and training variability}
% src: bib:muller22_in_the_wild
Cross-corpus evaluation tests more than the ability to reject a new synthesis algorithm. In-the-Wild collects
authentic and manipulated speech of public figures from online media, and demonstrates poor generalization of
ASVspoof-trained detectors to that setting \citep{muller22_in_the_wild}. Recording conditions, speech content,
speakers, and generation pipelines can change together. A target EER summarizes the resulting discrimination,
whereas the transferred FRR and FAR additionally reveal which class a deployed rule harms.

% src: bib:muller21_silence
Dataset cues can make this distinction harder to interpret. \citet{muller21_silence} identify label-correlated
leading and trailing silence in ASVspoof data; silence-duration classifiers are predictive, and training established
CMs after trimming silence changes their detection performance. Their evidence motivates interventions on the
training audio, rather than treating a test-time trim alone as evidence that a detector has learned a particular
causal mechanism. Our paired training design asks whether removing this cue changes the transferred operating
behavior, with source performance serving as a manipulation check.

% src: bib:wang25_asvspoof5_db
ASVspoof~5 explicitly addresses collection artefacts. Its shortcut-reduction pipeline normalizes peak amplitude,
randomly trims leading and trailing non-speech, and varies the retained utterance duration. The database paper
states that this pipeline is applied to the development and evaluation sets, but not the training set
\citep{wang25_asvspoof5_db}. Consequently, a detector trained on that release is not automatically trained on
shortcut-reduced audio. Training variability is another source of uncertainty:
\citet{wang21_comparative} show that changing initialization can produce statistically distinguishable performance
for the same CM recipe. Replication over training runs therefore complements comparisons of released checkpoints.

\subsection{Countermeasure architectures and probability calibration}
% src: bib:tak21_rawnet2
The checkpoints in this audit span distinct representation choices. RawNet2 learns directly from waveforms
\citep{tak21_rawnet2}, while Res-TSSDNet uses a residual time-domain network \citep{hua21_res_tssdnet}.
AASIST integrates spectral and temporal graph-attention representations, and AASIST-L is its lighter variant
\citep{jung22_aasist}. These models support an assessment across architectures without attributing every difference
in transfer behavior to a single back end.

% src: bib:tak22_xlsr_aasist
Self-supervised speech representations broaden the feature space available to CMs. \citet{tak22_xlsr_aasist}
combine a wav2vec~2.0 front end with an AASIST back end and data augmentation. XLS-R supplies multilingual pretrained
speech representations \citep{babu21_xlsr}; XLSR-Mamba combines that representation family with a dual-column
bidirectional state-space back end \citep{xiao25_xlsrmamba}. Improvements in discrimination under these designs do
not themselves establish that a source decision threshold will transfer.

% src: bib:pascu24_calibrated
\citet{pascu24_calibrated} address both generalization and prediction reliability using frozen self-supervised
representations with logistic regression. They assess discrimination with EER and probability calibration with
expected calibration error. ECE describes agreement between predicted probabilities and observed class frequencies
within probability bins; it does not evaluate a particular source threshold after transfer. Including a detector designed
with this calibration objective allows the audit to examine that distinction empirically, without presuming that
probability calibration removes operating-point drift.

\subsection{Channel, codec, and multilingual shift}
% src: bib:liu23_asvspoof2021
Communication processing can alter detection cues in both genuine and manipulated speech. ASVspoof~2021 DF
includes media compression, while ASVspoof~5 applies encoding and compression, including EnCodec, to evaluation
utterances from both classes \citep{liu23_asvspoof2021,wang25_asvspoof5_db}. The communication benchmark of
\citet{realworld_comms24} investigates combinations of codecs and packet loss. These studies motivate controlling
channel changes when comparing clean and perturbed scores, and reporting the class-specific errors of the
transferred rule in addition to threshold-optimized detection performance.

% src: bib:wu24_codecfake
Neural codecs also participate in speech generation. CodecFake trains countermeasures on codec-resynthesized
speech to improve detection of codec-based synthesis \citep{wu24_codecfake}.
\citet{xiao26_resynth} examine the resulting label ambiguity: a codec can be either a transmission mechanism or
part of a synthesis pipeline. The codec operation alone therefore does not specify the audit's target label.
Our channel stress tests preserve the clean utterance's label and state that assumption explicitly; this assesses
the treatment of genuine speech carried through a channel, rather than defining every resynthesis as an attack.

% src: bib:mlaad24
MLAAD expands evaluation beyond predominantly English benchmarks by providing synthetic speech across languages
and synthesis architectures \citep{mlaad24}. Its distributed audio is synthetic; corresponding original human
recordings are supplied by the M-AILABS source corpus \citep{solak17_mailabs}. This distinction matters for threshold auditing: spoof-only
data can measure attack acceptance at a frozen threshold, but cannot identify false rejection or a complete HTER.
Language and synthesis coverage should therefore be reported separately from evidence about the operating rule on
bona fide traffic. The audit fixes the dataset version and pairing protocol in its study design.

\subsection{Selective decisions and transferred review policies}
% src: bib:geifman17_selective
Selective classification adds an abstention option, trading automatic-decision coverage for lower error among
the retained decisions \citep{geifman17_selective}. In speech screening, the abstention can become a review queue,
with acceptance and rejection on either side of a score band. A band chosen on development data remains a
transferred policy: its target review share and residual errors depend on the target score distribution. Our
three-zone assessment consequently measures coverage and class-specific automatic errors together. It does not
treat a source review budget as a guarantee of target workload, nor assume that a target-fitted oracle policy is
available during deployment.
